% ============================================================
% SECTION 11 — ALEMBIC (SCHEMA MIGRATIONS)
% ============================================================

\section{Schema Migrations (Alembic)}

\begin{tcolorbox}[summarybox={\iconGear[1.5] Safely Evolving the Schema}]
\color{textdark}
In the early days of development, RecoverAI V2 relied on SQLAlchemy's \code{Base.metadata.create\_all(bind=engine)} to create tables on startup. While convenient, this approach is destructive for production: it cannot alter existing tables, add columns, or apply index changes without dropping the entire database.

To solve this, the system uses \textbf{Alembic}, the industry-standard migration tool for SQLAlchemy. Alembic tracks database schema changes as versioned Python scripts, allowing the database to evolve safely alongside the application code.
\end{tcolorbox}

\vspace{0.8cm}

% --- HOW IT WORKS ---
\subsection{How Alembic Integrates}

Alembic operates as a separate CLI tool (\code{alembic upgrade head}) that runs before the FastAPI application boots. 

\begin{center}
\begin{tikzpicture}[
    node distance=1.5cm,
    script/.style={rectangle, draw=secondary, thick, fill=bglight, minimum width=3cm, minimum height=1cm, rounded corners=4pt, font=\small\bfseries, drop shadow=black!5, align=center},
    db/.style={cylinder, draw=primary, thick, fill=primary!10, shape border rotate=90, aspect=0.25, minimum width=2.5cm, minimum height=1.5cm, font=\small\bfseries, drop shadow=black!5, align=center},
    table/.style={rectangle, draw=textdark, dashed, fill=white, minimum width=2.5cm, minimum height=0.6cm, font=\scriptsize\ttfamily, align=center},
    arrow/.style={-{Stealth[scale=1.1]}, thick, draw=textdark}
]

    \node[script] (m1) {Migration 1\\(Baseline)};
    \node[script, below=0.5cm of m1] (m2) {Migration 2\\(Webhook Retry)};
    
    \node[db, right=3cm of m1, yshift=-0.5cm] (pg) {Database};
    
    \node[table, below=0.1cm of pg] (alembic) {alembic\_version};
    
    \draw[arrow] (m1.east) -- node[above, font=\scriptsize] {Applies schema} (pg.160);
    \draw[arrow] (m2.east) -- node[below, font=\scriptsize] {Alters table} (pg.200);
    
    \draw[->, thick, draw=secondary, dashed] (alembic.west) -- ++(-1,0) |- node[near start, below, font=\scriptsize, text=secondary] {Checks current state} (m2.west);

\end{tikzpicture}
\end{center}

\begin{tcolorbox}[infobox={The \code{alembic\_version} Table}]
When Alembic runs, it creates a single-row table in the database called \code{alembic\_version}. This table stores the hash of the most recently applied migration script. When a deployment occurs, Alembic reads this hash, compares it to the scripts in the repository, and applies only the missing updates in chronological order.
\end{tcolorbox}

\vspace{0.8cm}

% --- MIGRATION HISTORY ---
\subsection{Migration History}

The repository currently contains two migration files in the \filepath{backend/alembic/versions/} directory.

\subsubsection{1. The Baseline Migration}
\textbf{ID:} \code{647f86570217} | \textbf{File:} \filepath{647f86570217\_baseline\_schema.py}

This migration represents the consolidated state of the database after Batches 1 through 4.6. Instead of having dozens of tiny development migrations, the engineering team "squashed" the schema into a single baseline before preparing for production. It creates:
\begin{itemize}
    \item \code{transactions}
    \item \code{recovery\_attempts}
    \item \code{audit\_logs}
    \item \code{idempotency\_records}
    \item \code{webhook\_events} (Original schema)
\end{itemize}

\vspace{0.4cm}

\subsubsection{2. The Webhook Retry Remediation}
\textbf{ID:} \code{faecf6256136} | \textbf{File:} \filepath{faecf6256136\_add\_webhook\_retry\_fields.py}

Generated during Batch 5.3.1 to resolve the P1 Webhook Retry Starvation vulnerability discovered in the adversarial audit.

\begin{tcolorbox}[codebox={Batch 5.3.1 Migration Code}]
\begin{lstlisting}[language=Python]
def upgrade() -> None:
    # Add retry_count with a default of 0
    op.add_column('webhook_events', 
        sa.Column('retry_count', sa.Integer(), server_default='0', nullable=False)
    )
    # Add last_attempt_at for bounded sweep scheduling
    op.add_column('webhook_events', 
        sa.Column('last_attempt_at', sa.DateTime(timezone=True), nullable=True)
    )

def downgrade() -> None:
    op.drop_column('webhook_events', 'last_attempt_at')
    op.drop_column('webhook_events', 'retry_count')
\end{lstlisting}
\end{tcolorbox}

\vspace{0.8cm}

% --- BEST PRACTICES FOR FUTURE DEVELOPMENT ---
\subsection{Best Practices for Future Changes}

As RecoverAI V2 moves toward production, any future changes to the database schema must strictly adhere to the following rules:

\begin{tcolorbox}[warningbox={Production Schema Rules}]
\begin{enumerate}
    \item \textbf{No Model Edits Without Migrations:} Changing \filepath{app/models/db\_models.py} without generating an Alembic migration will result in a fatal schema mismatch in production.
    \item \textbf{Autogenerate Support:} Use \code{alembic revision --autogenerate -m "description"} to automatically generate migration scripts based on changes to the SQLAlchemy models.
    \item \textbf{Manual Review Required:} Autogenerated scripts must be manually reviewed. Alembic often struggles to correctly infer table renames or complex constraint changes.
    \item \textbf{SQLite Limitations:} SQLite does not fully support the \code{ALTER TABLE} command (e.g., dropping columns is highly problematic in SQLite but trivial in PostgreSQL). Alembic handles this via "batch mode" recreations, but testing migrations on PostgreSQL is mandatory.
\end{enumerate}
\end{tcolorbox}
